\section{Related Work}

Several studies use SIH/SUS data to study where Brazilians are hospitalized.
Oliveira et al.~\cite{Oliveira2004TerritoriosSistema} mapped hospital care networks and found that most admissions still happen in the patient's own municipality, while the remaining flows point toward regional hubs.
Rocha et al.~\cite{Rocha2015HospitalizationFlow} measured patient migration between health regions in S\~{a}o Paulo and classified each region as attracting or losing patients.
Xavier et al.~\cite{Xavier2019RegioesSaude} used the origin and destination of admissions to build health regions and to locate care gaps.
These works report indices, tables, and static maps, but they rarely let the reader explore the flows directly.

Information visualization has long studied how to draw origin-destination (OD) data.
Drawing every flow as a straight line produces heavy occlusion, so much of the work looks for ways to reduce clutter.
Wood et al.~\cite{Wood2010VisualisationOrigins} replaced lines with small gridded maps (OD maps) to avoid crossing flows, while Jenny et al.~\cite{Jenny2018DesignPrinciples} proposed design principles for line-based flow maps and Graser et al.~\cite{Graser2019Untangling} compared techniques for untangling OD flows in a GIS.
For data that changes over time, Flowstrates~\cite{Boyandin2011Flowstrates} separates origins, destinations, and a temporal heatmap.
We borrow the clutter-reduction goal from this line of work and apply it through three concrete choices: curved arcs in place of straight lines, aggregation by administrative level so that only a tractable number of arcs is drawn at once, and a direction filter that hides the arcs the user is not inspecting.

Interaction is what lets users work with large and multivariate data.
Shneiderman's mantra of overview first, zoom and filter, then details on demand~\cite{Shneiderman1996EyesHaveIt}, together with the interaction categories of Yi et al.~\cite{Yi2007TowardDeeper}, describes the operations our tool depends on.
The three-level drill-down and the specialty and time filters put these ideas to work on the hospitalization data.

General flow-map tools such as kepler.gl and flowmap.blue can plot OD data, but they are generic plotters that a user must configure for each dataset.
Our work connects the two threads above and adds three features that a generic plotter does not provide on this data.
First, a three-level drill-down bound to the Brazilian administrative hierarchy, from federative unit to municipality.
Second, a precomputed aggregation that keeps the full 183 million record dataset interactive.
Third, panels built for the access questions at hand: an attraction ranking and a resident-versus-incoming seasonal comparison.
